% Иллюстрация элементарного телесного угла:
% площадка на единичной сфере между theta, theta+dtheta и phi, phi+dphi
% со сторонами dtheta и sin(theta) dphi.
% Компиляция: pdflatex solid-angle-figure.tex
\documentclass[tikz,border=8pt]{standalone}
\usepackage[T2A]{fontenc}
\usepackage[utf8]{inputenc}
\usepackage[russian]{babel}
\usetikzlibrary{arrows.meta,calc}

\begin{document}
\begin{tikzpicture}[
    >={Stealth[length=2.4mm]},
    axis/.style   = {->, semithick},
    constr/.style = {thin, gray!55},
    patch/.style  = {semithick, red!70!black, fill=red!18},
  ]
  % проекция: sx = 2.6 sin(t)sin(p);  sy = 2.6(0.971 cos(t) - 0.238 sin(t)cos(p))

  % ================= сфера =================
  \draw[semithick] (0,0) circle (2.6);
  % экватор: передняя половина сплошная, задняя пунктиром
  \draw[thin, domain=-90:90, samples=60]
        plot ({2.6*sin(\x)}, {-0.619*cos(\x)});
  \draw[thin, dashed, domain=90:270, samples=60]
        plot ({2.6*sin(\x)}, {-0.619*cos(\x)});
  % полярная ось
  \draw[dashed, thin] (0,0) -- (0,2.5246);
  \draw[axis] (0,2.5246) -- (0,3.35) node[above] {$z$};
  \node[gray!30!black] at (-0.24,-0.18) {$O$};

  % ================= построения: меридианы и параллели =================
  % меридианы phi = 20 и phi = 48
  \draw[constr, domain=0:90, samples=50]
        plot ({0.8892*sin(\x)}, {2.5246*cos(\x)-0.5816*sin(\x)});
  \draw[constr, domain=0:90, samples=50]
        plot ({1.9318*sin(\x)}, {2.5246*cos(\x)-0.4139*sin(\x)});
  % параллели theta = 50 и theta = 68
  \draw[constr, domain=14:85, samples=60]
        plot ({1.992*sin(\x)}, {1.6234-0.4740*cos(\x)});
  \draw[constr, domain=14:85, samples=60]
        plot ({2.410*sin(\x)}, {0.9467-0.5736*cos(\x)});

  % ================= площадка dOmega =================
  \draw[patch] (0.681,1.178) -- (1.480,1.306) -- (1.791,0.563) -- (0.824,0.407) -- cycle;

  % радиус к площадке и угол theta
  \draw[thin] (0,0) -- (0.681,1.178);
  \draw[gray!30!black] (0,0) ++(90:0.55) arc[start angle=90, end angle=60, radius=0.55];
  \node[gray!30!black] at (0.24,0.98) {$\theta$};

  % dphi на экваторе
  \draw[red!70!black, semithick, domain=20:48, samples=20]
        plot ({2.6*sin(\x)}, {-0.619*cos(\x)});
  \node[red!70!black] at (1.47,-0.78) {$\mathrm{d}\varphi$};

  % ================= подписи сторон =================
  \node[red!70!black] at (2.02,0.95) {$\mathrm{d}\theta$};
  \node[red!70!black, anchor=west] at (2.78,0.20) {$\sin\theta\,\mathrm{d}\varphi$};
  \draw[gray!70, thin] (2.72,0.21) -- (1.55,0.51);

\end{tikzpicture}
\end{document}
